\section{Modularity with Involutive Algebras}
\applabel{involutive}
We recount Jacobs's account, albeit in a more advanced categorical jargon, and
use it to prove a generic representation theorem for involutive frals
and frexes. We don't use this development elsewhere in this
article.

\begin{wrapfigure}{r}{.2\textwidth}
  \vspace{-.5\baselineskip}
  \includegraphics{diagram-08.mps}
  \vspace{-1.5\baselineskip}
\end{wrapfigure}
Jacobs appeals to the Baez-Dolan \emph{microcosm
  principle}~\cite{baez-dolan:microcosm} --- an algebraic structure on
an object
requires a compatible structure on its category of context---and
defines the following concepts. An \emph{involutive structure} on a
category $\C$ is a pair $\pair{\involve{(-)}}{\invlaw}$ consisting of a
functor $\involve{(-)}:\C \to \C$, called the \emph{involution}, and a
natural isomorphism $\invlaw : \iinvolve{(-)} \to -$, called the
\emph{involution law}, satisfying the condition on the right.
This definition is equivalent to Jacbos's, but reverses the direction of the involution
law.

For example, each category has an involutive structure given by the
identity functor as involution and the identity natural transformation
as the involutive law. This structure, which we call the
\emph{trivial} involutive structure, may seem degenerate, but it plays
an important role in our development.

The motivating example is $\Monoid{}$, the category of monoids. It has
the following non-trivial involutive structure.  Given a monoid \avar,
construct another monoid $\involve\avar$ with the operation reversed:
$\involve\avar\sem{\cdot}(x,y)\definedby \avar\sem\cdot(y,x)$. If
$\hName{} : \avar{} \to \bvar{}$ is a monoid homomorphism, then the
same underlying function provides a monoid homomorphism $\overline
\hName{} : \involve{\avar{}} \to \involve{\bvar{}}$. These maps define
an involution functor $\involve{(-)} : \Monoid \to \Monoid$. The
identity function is then a monoid isomorphism $\invlaw \definedby
(\fun xx) : \iinvolve\avar \to \avar$, the required involution law. We
have similar involutive structures on other categories, given by
ordinary or commutative: semi-groups, monoids, groups,
semirings and rings, and so on.

\begin{wrapfigure}{r}{.12\textwidth}
  \vspace{-1.4\baselineskip}
  \includegraphics{diagram-07.mps}\label{triangle-page}
  \vspace{-2\baselineskip}
\end{wrapfigure}
To see the microcosm principle in action, note that a function
$\hName{} : \U\,\avar \to \U\,\avar$ makes a monoid $\avar$ into an
involutive monoid if and only if (1) it is a monoid homomorphism $\hName{}
: \involve\avar \to \avar$, so $\hName(\x\cdot \y) =
\hName\,\y\cdot\hName\,\x$, and (2) the diagram on the right commutes, so
$\hName(\hName\,\x) = \x$.  These two conditions categorify the notion
of an involutive monoid, so we can define it in any involutive
category, not just $\Monoid$. Jacobs calls these \emph{self-conjugate}
objects, and we will study them in more detail soon.

\begin{wrapfigure}{r}{.21\textwidth}
  \vspace{-1.1\baselineskip}
  \includegraphics{diagram-09.mps}
  \vspace{-2\baselineskip}
\end{wrapfigure}
Packaging this structure, an \emph{involutive category} $\C =
\triple{\cat\C}{\involve{(-)}}{\invlaw}$ is an
\IdrisFunction{o}rdinary category $\cat\C$ equipped with an involutive
structure $\pair{\involve{(-)}}{\invlaw}$. An \emph{involutive
  functor} ${F : \C[2] \to \C[1]}$ between involutive categories is a
pair $\pair{\cat F}{\dist[F]}$ consisting of an ordinary functor $\cat
F : \cat {\C[2]} \to \cat \C$ and a natural transformation $\dist[F]
:\cat F\involve{(-)} \to \involve{\cat F(-)}$ called its
\emph{distributive law}, satisfying the compatibility condition on the
right. Such distributive laws are natural isomorphisms.

The canonical example is the forgetful functor $\U : \Model \Pres \to
\Set$ from the category of models of some presentation $\Pres$ to the
category of sets and functions. This functor has an involutive functor
structure w.r.t.~an involutive structure on $\Model\Pres$, when the
involution of an algebra only changes the operations of the algebra,
but not its carrier. Note the role that the
trivial involutive structure on $\Set$ plays. All the examples above
of monoid varieties and the semi-ring
varieties w.r.t.~the operation-reversal and trivial involutive
structures
have such involutive forgetful functors.

\begin{wrapfigure}{r}{.16\textwidth}
  \vspace{-1.1\baselineskip}
  \includegraphics{diagram-10.mps}
  \vspace{-2\baselineskip}
\end{wrapfigure}
An \emph{involutive natural transformation} $\alpha : F \to G$ between
two involutive functors is an ordinary natural transformation $\alpha
: \cat F \to \cat G$ between their underlying ordinary functors that
moreover satisfies the condition on the right. As Jacobs comments, we
therefore have a 2-category $\ICAT$ consisting of involutive
categories, functors, and natural transformations, and we may derive
involutive adjunctions as two involutive functors and two involutive
natural transformations satisfying the triangle laws.

We can turn an ordinary adjunction
into an involutive one when one of the functors is involutive:
\begin{proposition}
  Let $G : \C[3] \to \C$ be an involutive functor, and $\cat F
  \leftadjointto \cat G$ be a left-adjoint to the ordinary functor
  underlying $G$ with unit $\unit$ and counit $\counit$.
%
  Set $\dist[F]_{\X} : \cat F \involve\X \to \involve{\cat F\X}$ to be
  the mate of the composite
  \(
  \involve\X \xto{\involve\unit} \overline{\cat G\cat F \X}
  \xto{(\dist[G])^{-1}} \cat G\involve{\cat F\X}
  \).
  Then (1) $\dist[F]$ equips $\cat F$ with an involutive functor structure
    $F : \C \to \C[3]$; and (2) $F \leftadjointto G$ is an involutive adjunction with unit
    $\unit$ and counit $\counit$.
\end{proposition}

As a consequence, the free model functors for models in which the
forgetful functor is involutive are all involutive adjunctions. This
consequence covers our monoid and semi-ring varieties of interest,
namely ordinary and commutative semi-groups, monoids, groups,
semi-rings and rings with or without a unit. The distributive laws in
these examples are given by the mate of the function:
\(
\fun x\unit x : \involve\X = \X \xto{\involve\unit = \unit } UF\X
\xto{\dist[U] = \fun tt} U\involve{F\X}
\). %
%
One might be tempted to think that the resulting distributive law is
the identity homomorphism, because the mate of the unit of an
adjunction is the identity function. It is not the case. When we take
the mate, we take into account the algebra structure of
$\involve{F\X}$, which may change the interpretation of the
operations, and consequently changes the resulting mate homomorphism.
For the non-trivial involutive structures over monoid and semi-ring
varieties, the distributive law will reverse the relevant binary
operation.


\begin{wrapfigure}{r}{.12\textwidth}
  \vspace{-1\baselineskip}
  \includegraphics{diagram-11.mps}
  \vspace{-1\baselineskip}
\end{wrapfigure}
A \emph{self-conjugate} object $\avar{}$ in an involutive category
$\C[3]$ is a pair $\pair{\obj \avar}{\aconj[\avar]}$ consisting of an
object $\obj \avar$ in $\C[3]$, and an $\C[3]$-morphism $\aconj[A] :
\involve {\obj\avar} \to \obj\avar$, satisfying the triangle on the
right. As we saw on p.~\pageref{triangle-page}, self-conjugate monoids
are involutive monoids, and more generally, self-conjugate
semi-groups, groups, semi-rings, rings, etc.~are the involutive
ones. A \emph{homomorphism} $\hName : \avar \to \bvar$ of
self-conjugate objects is a homomorphism\\[-\baselineskip]

\begin{wrapfigure}{l}{.12\textwidth}
  \vspace{-1\baselineskip}
  \includegraphics{diagram-12.mps}
  \vspace{-2\baselineskip}
\end{wrapfigure}\noindent{}%
$\hName : \obj\avar \to \obj\bvar$ between their underlying objects
that moreover satisfies the condition on the left. This condition
generalises the usual condition of involutive monoid homomorphisms and
so on. Since homomorphisms of self-conjugate objects compose and
contain the identities, they form a category which we denote by
$\SC\C[3]$. Jacobs shows that the forgetful functor $\U : \SC\Set \to
\Set$ has a left adjoint $\FSC : \Set \to \SC\Set$ sending each set
$\X$ to the coproduct of two copies of itself, i.e.~by tagging each
element with a boolean, and the self-conjugation structure flips this boolean
$\FSC\X \definedby \pair{\BoolX}{\fun {(\bvar,\x)}(\lnot\bvar, \x)}$.

It will pay-off immediately to include one more level of abstraction.
Jacobs (Lemma~6) shows that the $\SC$-construction extends to a
2-functor $\SC[] : \ICAT \to \ICAT$. We recall the remaining
structure. The action on objects of $\ICAT$ equips the category
$\SC\C[3]$ with an involutive structure, sending each self-conjugate
object $\avar$ to the self-conjugate object $\involve\avar \definedby
\pair{\involve{\obj\avar}}{\involve{\aconj[\avar]} :
  \iinvolve{\obj\avar} \to \involve{\obj\avar}}$. The action on the%
\\[-\baselineskip]

\begin{wrapfigure}{l}{.18\textwidth}
  \vspace{-1\baselineskip}
  \includegraphics{diagram-13.mps}
  \vspace{-1\baselineskip}
\end{wrapfigure}\noindent{}%
morphisms of $\ICAT$, sends an involutive functor
$F : \C[2] \to \C$
to the involutive functor $\SC F : \SC\C[2] \to \SC\C$ mapping each
self-conjugate object $\avar$ to the self-conjugate object
$\pair{\cat
  F\obj \avar}{\aconj[\SC F \avar] : \involve{\cat F\obj
    \avar}\xto{(\dist[F])^{-1}} \cat F\involve{\obj\avar} \xto{\cat
    F\aconj[\avar]}}$, and acting as $\cat F$ on self-conjugate
homomorphisms. The action on 2-cells sends each involutive natural
transformations to itself, i.e.~a
natural transformation between
involutive functors also preserves the resulting self-conjugated
structure. The
forgetful functor $\U : \SC\C[3] \to \C[3]$ is then natural as on the
left.

\begin{wrapfigure}{r}{.22\textwidth}
  \vspace{-1\baselineskip}
  \includegraphics{diagram-14.mps}
  \vspace{-2\baselineskip}
\end{wrapfigure}
We profit off of this obscene level of abstraction immediately:
2-functors preserve all 2-adjunctions, since they transport the
triangle equalities to the appropriate triangle equalities. Therefore,
if we have an involutive adjunction $F \leftadjointto G : \C[3] \to
\C$ where $\C$ has a free self-conjugate object adjunction $\FSC^{\C} \leftadjointto \U
: \SC \C \to \C$, we get the free self-conjugate $\C[3]$-object on $\X$
as the composite $F(\FSC^{\C} \X)$ completely structurally, as on the
right. Applying this result to frals, we get the following
generalisation of \propref{involution}(fral):
\begin{proposition}\label{prop:involutive algebra fral}
Let $\Pres$ be a presentation equipped with an involutive structure
over $\Model\Pres$ and an involutive forgetful functor structure. Let $\pair\A\Env$ be any free $\Pres$ model over
the product $\BoolX$. Then the following structure exhibits $\A$ as
the free self-conjugate $\Pres$-model over $\X$:
\[
\begin{array}[b]{@{}l@{}}
\aconj[\A] : \involve{\A} \xto{\dist[-1]} \A \xto{\bind(\lnot \times \id)} \A{}
\\
\Env' : \X \xto{\fun \x\spair\IFalse\x} \BoolX
\end{array}
\qquad
\begin{array}[b]{@{}l@{}}
(\bind'\f) : \A \xto{\bind\fun{\spair\bvar\x}\begin{cases}
    b = \ITrue: & \aconj[\avar](\f\,\x) \\
    b = \IFalse: & \f\,\x
  \end{cases}
} \avar
\end{array}
\]
\end{proposition}

Having dealt with the fral, we turn to the frex. Jacobs proves that if
$\triple{\avar_1 + \avar_2}{\injection_1}{\injection_2}$ is a
coproduct of $\avar_1$ and $\avar_2$ in an involutive category, then
$\triple{\involve{\avar_1 +
    \avar_2}}{\involve{\injection_1}}{\involve{\injection_2}}$ is a
coproduct of $\involve{\avar_1}$ and $\involve{\avar_2}$. The unique
cotupling morphism in the universal property, for each $\hName_i : \involve{\avar_i} \to \bvar$, is
\(
\involve[\hName_1, \hName_2\involve] : \involve{\avar_1 + \avar_2}
\xto{[\involve\hName_1 \compose \invlaw^{-1},
    \involve\hName_2\compose\invlaw^{-1}]} \iinvolve \bvar \xto{\invlaw} \bvar
\).
If each $\avar_i$ has a self-conjugate structure $\aconj[i] :
\involve{\avar_i} \to \avar_i$, then the coproduct $\avar_1 + \avar_2$
has a self-conjugate structure given by $\aconj[1] + \aconj[2]
\definedby \involve[\injection_1\compose \aconj[1],
  \injection_2\compose \aconj[2]\involve] : \involve{\avar_1 +
  \avar_2} \to \avar_1 + \avar_2$.
Since the frex $\avar{}[\X]$ can be construct as the coproduct of the
model $\avar$ with the fral on $\X$, we generalise
\propref{involution}(frex):
\begin{proposition}\label{prop:involutive algebra frex}
Let $\Pres$ be a presentation equipped with an involutive structure
over\/ $\Model\Pres$ and an involutive forgetful functor structure, and $\avar$ be a self-conjugate $\Pres$-model. Let
$\triple\A\Var\Embed$ be any $\Pres$-frex of $\obj\avar$ by the
product $\BoolX$. Then the following structure exhibits $\A$ as the
frex of the self-conjugate $\Pres$-model \avar{} by $\X$, for $\hName :
\avar \to \cvar$ involutive monoid homomorphism and function $\env :
\X \to \cvar$:
\begin{gather*}
\aconj : \involve \A \xto{\bracks{\obj\avar\xto{\invlaw^{-1}} \iinvolve{\obj\avar} \xto{\involve {\aconj[\avar]}} \involve\avar \xto{\involve\Embed} \involve\A
, \BoolX \xto{\lnot \times \id} \BoolX \xto{\Var} \U\,\A \xto{\dist[-1]} \U\,{\involve\A}}} \iinvolve\A \xto{\invlaw} \A\\
\Var' : \X \xto{\fun \x\spair\IFalse\x} \BoolX \xto{\Var} \A
\qquad
\Embed : \obj\avar \xto{\Embed} \A\\
[\hName,\env] : \A \xto{\bracks{h, \fun{\spair \bvar x}\begin{cases}
           \bvar = \ITrue: & \aconj[\cvar](\env\,\x) \\
           \bvar = \IFalse: & (\env\,\x)
         \end{cases}}
}
\cvar
\end{gather*}
\end{proposition}

\Frexlib{} does not yet implement this proposition in its full
generality, since it would require a substantial amount of
additional infrastructure, either inside \Frexlib{} or as part of a
category-theory library for \idris{}. For example, the type of the
construction requires a categorical equivalence
between some $\Model\Pres'$ for the presentation $\Pres'$ of involutive
$\Pres$-models and the self-conjugate $\Pres$-models.
\Frexlib{} currently only implements the
special case of \propref{involution}, with its specialised proof.
